\label{chap:introduction}

Recent technological advances have led manufacturing companies to equip their machines with increasingly sophisticated sensors and interaction interfaces. At the same time, the widespread adoption of Artificial Intelligence (AI) technologies, particularly Neural Networks (NNs), has encouraged companies to investigate how these methods could improve both machine functionalities and operator interfaces.

As a result, these technological developments have enabled not only the creation of more advanced machinery, but also the development of Intelligent Decision Support Systems (IDSS)~\cite{phillips2013intelligent} capable of assisting machine operators during manufacturing activities.

This doctoral thesis was developed in collaboration with SCM Group S.p.A., an Italian company and world leader in the production of machinery for processing a wide range of materials, including wood, plastic, glass, stone, metal, and composite materials. Its machines serve several industrial sectors and support different manufacturing processes. Consequently, the machinery produced by SCM Group varies significantly in functionality and features. Despite these differences, all machines involve interaction with a human operator during the initial and final stages of processing and the need to satisfy operational and safety constraints throughout the manufacturing process.

Specifically, during the initial stage of production, the machine operator configures the parameters governing the manufacturing process, while in the final stage the quality of the finished product is evaluated. The machines may include several configurable and controllable components whose operation must comply with operational and safety constraints aimed at protecting both the machinery and the operator.

The physical processes underlying manufacturing operations directly influence the quality of the final product. However, during machine operation, operators often lack the time or expertise required to analyze these processes in detail. Consequently, they rely on the machine interface to provide guidance and support for decision making.

The objective of this work is to improve human-machine interaction developing intelligent decision support system aimed at optimizing tasks performed by the operator during the initial stages of machine configuration.

The remainder of this thesis investigates two optimization problems. The first problem, referred to as \textit{Fixture Layout Optimization}, addresses the positioning of fixtures used to support wooden workpieces during machining, taking as reference a \textit{work center} for wood processing. The second problem, referred to as \textit{Heating Process Optimization}, concerns the determination of machine parameters required to achieve a target temperature distribution within the core of a plastic sheet, taking as reference a \textit{thermoforming} machine.

The thesis is therefore divided into two main parts. The first part presents the Fixture Layout Optimization problem, along with the techniques and tools adopted to address it, their limitations, and the novel solutions proposed to overcome them. The second part focuses on the Heating Process Optimization problem, describing the data-driven approach adopted and the optimization techniques used to solve it.
